% !TEX root = ../jasonliu2026_Thesis.tex

\setcounter{chapter}{-1}

\chapter{Preamble}

\section*{Abstract}
\addcontentsline{toc}{section}{Abstract}

Automatically generated unit tests suffer from uninformative identifier names such as \texttt{var0} and \texttt{test42}, which impede comprehension and reduce the practical value of the generated output. 
This paper presents an LLM-based renaming pipeline that renames both local variables and test method names in generated Java tests by prompting a locally served open-source model with structured context extracted from the test body. We evaluate the pipeline using  Qwen2.5-Coder 7B across six fine-tuning checkpoints and GPT-5 mini as a zero-shot closed-source baseline, compared against the non-LLM baselines GGNN and RefBERT on a 100-sample subset of the \texttt{methods2test} dataset. We further introduce CodeReader, a configurable open-source readability grading tool that scores renamed output using a weighted ensemble of LLM judges against an explicit rubric, enabling reproducible evaluation without requiring a human study. Both LLM models far surpass the non-LLM baselines (F1\,=\,0.760 and 0.756 vs.\ 0.019 and 0.308). Fine-tuning improves F1 further to 0.839 while saturating after just 15\% of training steps. A counterintuitive finding emerges: fine-tuning increases token-level accuracy but reduces the LLM readability score, while GPT-5 mini achieves the highest readability score of all variants despite comparable F1, suggesting that token-level accuracy and  holistic readability measure fundamentally different properties. These results highlight the need to report both metrics when evaluating identifier renaming tools.

